\documentclass[12pt]{book}
\usepackage[utf8]{inputenc}
\usepackage[T1]{fontenc}
\usepackage{amsmath,amssymb,amsfonts}
\usepackage{amsthm}

\begin{document}

\setcounter{374}
\section{The duality theorem for proper morphisms.}

In this section we prove the long-awaited duality theorem for a proper morphism
\(f: X \to Y\) of locally noetherian preschemes.
We will suppose the existence of a residual complex \(K^{\cdot}\) on \(Y\).
Then \(K^{\cdot}\) and \(f^{\Delta} K^{\cdot}\) give rise to pointwise dualizing functors \(D_{Y}\) and \(D_{X}\) on \(Y\) and \(X\), respectively,
and we express the duality theorem as
\[D_{X}\left(F^{\cdot}\right) \stackrel{\sim}{\to} D_{Y}\left(\mathbf{R} f_{*} F^{\cdot}\right)\]
for \(F^{\cdot} \in D_{qc}(X)\).

Before proving the theorem, we must show that the trace \(Tr_{f}\) of [VI 4.2] agrees with \(Tr_f\) of [III 6.5]
in the case of a finite or projective morphism.

So let \(f: X \to Y\) be a finite morphism, and assume that
\(Y\) is noetherian and has finite Krull dimension.
Let \(K^{\cdot}\) be a residual complex on \(Y\).
Then we construct an isomorphism
\[\alpha_{f}: Q f_{*} f^{\Delta} K^{\cdot} \to \mathbf{R} f_{*} f^{b} Q K^{\cdot}\]
as follows:

\setcounter{page}{375}
\[Q f_{*} f^{\Delta} K^{\cdot} \stackrel{\xi_{f_{*}}}{\to} \mathbf{R} f_{*} Q f^{\Delta} K^{\cdot}
\stackrel{\psi_{f}}{\to} \mathbf{R} f_{*} f^{\nabla} K^{\cdot}\]
where as usual, \(Q\) denotes the natural map from complexes to elements of the derived category;
\(\xi_{f_{*}}\) is the isomorphism of [I.5.1];
\(\psi_{f}\) is the isomorphism of [VI 3.1c];
the map (I) is the definition of \(f^{Y}\) [VI],
and (2) is the isomorphism of [VI 1.1c]
(here is where we need the hypothesis that \(Y\) is noetherian of finite Krull dimension).

\begin{proposition}
Let \(f: X \to Y\) be a finite morphism, with \(Y\ noetherian of finite Krull dimension,
and let \(K^{\cdot}\) be a residual complex on \(Y\).
Then there is a commutative diagram
where \(Tr_{f}\) is the trace of [VI 4.2] and \(Tr_f\) is the trace of [III 6.5].
\end{proposition}

\setcounter{page}{376}
\begin{proof}
Follows immediately from TRA 2 [VI 4.2] and the definition of \(\rho_{f}\) [VI 4].
\end{proof}

Now let \(Y\) be a noetherian prescheme of finite Krull dimension,
and let \(f: X=\mathbb{P}_{Y}^{n} \to Y\) be the structural map of an \(n\)-dimensional projective space over \(Y\).
Let \(K^{\cdot}\) be a residual complex on \(Y\).
Then we define an isomorphism
\[\beta_{f}: Q f_{*} f^{\Delta} K^{\cdot} \stackrel{\sim}{\to} \mathbf{R} f_{*} f^{*} Q K^{\cdot}\]
similar to the map \(\alpha_{f}\) above, using \(\xi_{f_{*}}\),
the map \(\varphi_{f}\) of [VI 3.1d], the definition of \(f^{z}\) [VI], and [VI 1.1c].

\begin{proposition}
Let \(Y\) be a noetherian prescheme of finite Krull dimension,
let \(X=\mathbb{P}_{Y}^{n}\), let \(f: X \to Y\) be the projection,
and let \(K^{\cdot}\) be a residual complex on \(Y\).
Then there is a commutative diagram
where \(Tr_{f}\) is the trace of [VI 4.2], and \(Trp_{f}\) is the trace of [III 6.3].
\end{proposition}

\setcounter{page}{377}
\begin{proof}
Choose a section \(s: Y \to X\) of \(f\), and consider the following diagram:
where the notations have the sources indicated,
and the second vertical arrow is obtained by sandwiching \(\alpha_{s}\)
in the middle of the four isomorphisms which define \(\beta_{f}\).

Now we make the following observations:
i). The composition of the upper row is the identity on \(QK^{\cdot}\).
This follows from TRA 1 [VI 4.2].

2). The composition of the lower row is also the identity on \(QK^{\cdot}\).
This is the statement of [III 10.1].

3). The left-hand square is commutative. This follows from VAR 5 [VI 3.1] and the definition of the isomorphisms \(\alpha_{s}\), \(\beta_{f}\) and (IV) of [VI 3.2].

\setcounter{page}{378}
4). Proposition 3.1 above.

5). \(c_{s,f}\) and \(Trp_{f}\) are isomorphisms by construction, hence \(Trf_{s}\) is an isomorphism.

6). \(c_{s,f}\) is an isomorphism, hence \(Q Tr_{s}\) and \(Q Tr_{f}\) are isomorphisms.
(Note incidentally we have used Theorem 2.1 above that \(Tr_{f}\) is a morphism of complexes,
in order to consider \(Q Tr_{f}\) in the first place.)

We conclude finally that the right-hand square is commutative.
q.e.d.

Now we come to the duality theorem itself.
Let \(f: X \to Y\) be a proper morphism of noetherian preschemes of finite Krull dimension.
Let \(K^{\cdot}\) be a residual complex on \(Y\).
(Note that the existence of a residual complex imposes a slight restriction on the preschemes considered [VI 1.1] and [V].)
We denote by \(D_{Y}\) (resp. \(\underline{D}_{Y}\)) the functor \(\mathbf{R}\Hom_{Y}^{\cdot}(\cdot, Q K^{\cdot})\),
and by \(D_{X}\) (resp. \(\underline{D}_{X}\)) the functor \(\mathbf{R}\Hom_{X}^{\cdot}(\cdot, Q f^{\Delta} K^{\cdot})\).
Then composing the morphism of [II.5.5] with \(\xi f_{*}\) and \(Q Tr_{f}\), we obtain the duality morphism

\setcounter{page}{379}
\[\theta_{f}: \mathbf{R} f_{*} \underline{D}_{X}\left(F^{\cdot}\right) \to \underline{D}_{Y}\left(\mathbf{R} f_{*} F^{\cdot}\right)\]
for \(F^{\cdot} \in D^{-}(X)\).

Applying the functor \(\mathbf{R}\Gamma(Y,\cdot)\) to both sides, and using [II 5.2] and [II 5.3] we obtain a global duality morphism
\[\theta_{f}: D_{X}\left(F^{\cdot}\right) \to D_{Y}\left(\mathbf{R} f_{*} F^{\cdot}\right).\]

Taking the cohomology of this, we get morphisms
\[\theta_{f}^{i}: Ext_{X}^{i}\left(F^{\cdot}, Q f^{\Delta} K^{\cdot}\right) \to Ext_{Y}^{i}\left(\mathbf{R} f_{*} F^{\cdot}, Q K^{\cdot}\right).\]

\begin{theorem}[Duality Theorem]
Let \(f: X \to Y\) be a proper morphism of noetherian preschemes of finite Krull dimension,
and let \(K^{\cdot}\) be a residual complex on \(Y\).
Then the morphisms \(\underline{\theta}_{f}\), \(\theta_{f}\), and \(\theta_{f}^{i}\) defined above are isomorphisms
for all \(F^{\cdot} \in D_{qc}^{-}(X)\).
\end{theorem}

(Note that the hypothesis of finite Krull dimension is needed only for the definition of \(\beta_{f}\) (cf. [II 5.5]).
If we restrict to bounded complexes \(F^{\cdot} \in D_{qc}^{b}(X)\),
we can state the theorem assuming only that \(X\) and \(Y\) are locally noetherian,
and that the fibres of \(f\) are of bounded dimension. The proof is the same.)

\setcounter{380}
\begin{proof}
We proceed in several steps, eventually using Chow's lemma to reduce to the case of projective space which we know already.

a) Clearly it is sufficient to show that \(\underline{\theta}_{f}\) is an isomorphism.
The question is local on \(Y\), so we may assume \(Y\) is the spectrum of a local ring.
In particular, we may assume that \(Y\) is noetherian, affine, and of finite Krull dimension.
Using the lemma on way-out functors, we reduce to the case of a single quasi-coherent sheaf \(F\) on \(X\).

b) Any quasi-coherent sheaf is the direct limit of its coherent subsheaves.
In particular, it is a quotient of a direct sum of coherent sheaves.
Thus using the lemma on way-out functors again, we reduce to the case of a direct sum of coherent sheaves.
Now since \(\mathbf{R}\Hom\) transforms direct sums in the first variable to direct products,
we reduce to the case of a single coherent sheaf \(F\) on \(X\).

c) Now since \(Y\) is affine, and all the sheaves considered are quasi-coherent,
it is enough to show \(\theta_{f}\) is an isomorphism.

d) If \(f: X \to Y\) and \(g: Y \to Z\) are two proper morphisms,
and \(K_{Z}^{\cdot}\) a residual complex on \(Z\),
and if we take \(K^{\cdot}=g^{\Delta} K_{Z}^{\cdot}\), then for any \(F^{\cdot} \in D_{c}^{-}(X)\) we have a commutative diagram
\end{proof}

\end{document}